\begin{table*}[t]
\centering
\caption{Minimal natural-seawater exposure matrix for initial qualification of AM HEAs, structured on ASTM G52 \cite{astm_g52}, with crevice assemblies per ASTM G78 \cite{astm_g78}. All cells: triplicate specimens; two surface conditions (as-built, machined/polished); two build orientations; reference alloys (LPBF 316L, wrought 316L, NiAl bronze, 2205) in identical racks; at least two sites with different temperature/salinity/biota.}
\label{tab:exposure}
\footnotesize
\begin{tabularx}{\textwidth}{@{}lllXX@{}}
\toprule
Zone & Configuration & Durations & Flow condition & Primary endpoints \\
\midrule
Atmospheric / splash & Flat panels, 90$^{\circ}$ racks \cite{astm_g52} & 1, 6, 12, 24 mo & Wet--dry cycling, natural splash & Mass loss (G1/G31 \cite{astm_g1,astm_g31}), pit depth (G46 \cite{astm_g46}), rust/film analysis \\
Tidal & Panels on tidal rack \cite{astm_g52} & 1, 6, 12, 24 mo & Semidiurnal immersion cycles & Pit statistics, biofilm sampling, retained tensile properties \\
Full immersion (1--3 m) & Panels + crevice assemblies (G78 washers) + HEA--steel couples \cite{astm_g78,astm_g71} & 1, 6, 12, 24 mo & Quiescent to local tidal flow & Crevice attack rating, galvanic current log, EIS at retrieval intervals \\
Velocity loop (laboratory) & Pumped seawater rig & 3--6 mo & 0.5--3 m/s (pump/heat-exchanger service) & Erosion-corrosion, passive-film stability under flow \\
\bottomrule
\end{tabularx}

\vspace{4pt}
\footnotesize High-velocity service is outside G52's scope and must be added as a separate loop. Endpoints at every retrieval: mass loss, maximum pit depth and rating, crevice inspection, retained mechanical properties, and biofilm/under-film sampling (Table~\ref{tab:biofilm}).

\end{table*}
